\chapter{Data Collection}\label{ch:data}

\section{Introduction}
\textslot{what was recorded, how, and how it became training data; walks through Figure~\ref{fig:datapipe}.}

\diagram{datapipe}{The data pipeline, from teleoperation to a versioned dataset on the Hub.}{fig:datapipe}

\section{Tasks}\label{sec:tasks}

Two tasks were recorded, on three datasets (Table~\ref{tab:datasets}, Figure~\ref{fig:tasks}).
\begin{itemize}
\item \textbf{Cube in box.} A green cube is picked up and put into a box. It is done when the cube is inside the box
      and the gripper is open. The first recording (v0) was a smaller, initial dataset; the second (v1) repeats the
      task with more demonstrations and the final camera arrangement.
\item \textbf{Sorting.} A red or a blue cube, at a fixed position, is put into the square of its colour; the two
      coloured squares swap sides between layouts. It is done when the cube is fully inside the square of its colour
      and the gripper has let go of it. Sorting needs the policy to read the scene (which colour, which side), not
      only to repeat one motion, and is the main task of the DRTC experiments.
\end{itemize}

\begin{table}[htbp]
  \centering
  \small
  \gentable{tab_datasets}
  \caption{Datasets (pinned Hub revisions). Joint units: degrees, or LeRobot's normalised range ($-100$ to $100$) for
           the cube-in-box recordings.}\label{tab:datasets}
\end{table}

\begin{figure}[htbp]
  \centering
  \begin{subfigure}[t]{0.48\textwidth}
    \centering
    \photo[\textwidth]{task_sort_one.jpg}{The sorting task: the cube and the coloured squares, seen by the scene camera.}
    \caption{Sorting.}
  \end{subfigure}\hfill
  \begin{subfigure}[t]{0.48\textwidth}
    \centering
    \photo[\textwidth]{task_cubebox.jpg}{The cube-in-box task.}
    \caption{Cube in box.}
  \end{subfigure}
  \caption{The two tasks.}
  \label{fig:tasks}
\end{figure}

\diagram{layouts}{The task layouts, seen from above (schematic, not to scale).}{fig:layouts}

\section{Recording}\label{sec:recording}

Demonstrations are recorded by teleoperation: the operator moves the leader arm by hand and the follower copies it,
while both cameras and the follower's joint positions are recorded at 30~fps. One command runs the whole recording
session from the registry's rig: it opens the arms and cameras, counts episodes with a fixed episode time and reset
time, and afterwards checks the recording (dropped frames, camera rates, joint limits), trims it, pushes it to the Hub
and writes its pinned registry entry.

\paragraph{Held-out episodes on top.} A recording of $N+K$ episodes is planned from the start: $N$ to train on and
$K$ held out, chosen by a seeded random split and never used for training. The held-out episodes pick the best
checkpoint during training, give the offline action error (Section~\ref{sec:offline-design}), and are what the replay
bench plays (Section~\ref{sec:bench-design}). Sorting has $80+8$ episodes, cube in box v1 $100+10$.

\paragraph{Layouts.} The object positions are taped on the task board. The layouts rotate in a fixed, balanced order
over the episodes (for sorting: red or blue cube, times the two sides of the squares, four combinations), so every
combination is equally represented in training and in the held-out episodes.

\section{Trimming and editing}\label{sec:trimming}

A recorded episode starts and ends with the arm at rest while the operator gets ready and lets go. Trained on such
episodes, a policy learns that standing still is a valid action at the start: published SO-101 results show that
\enquote{never starting} is a major failure mode, which trimming the idle start and end removed almost
entirely~\cite{valloss2026}. Every dataset is therefore kept in two copies:
\begin{itemize}
\item the \textbf{raw} copy, the recording as made, which is the only copy ever edited; and
\item the \textbf{trimmed} copy, rebuilt from the raw one after every change, which is the one models train on.
\end{itemize}
The automatic trim measures the recording's own joint noise floor and cuts each episode's start and end where the arm
is still, without re-encoding the videos (each episode's video window moves). Editing tools on top let the operator
inspect every episode, view it, delete bad episodes (the rest are renumbered), set a manual trim, undo the last edit
and push. Manual trims are stored inside the raw copy, so they travel with it. Deleting episodes changes the seeded
held-out split, so all editing happens before training.

\section{Versions and splits}\label{sec:versions}

A dataset in the registry is pinned to one Hub revision, with its episode and frame counts and its joint units; a
training run refuses a local copy that does not match. The trimmed copy has the plain name and the raw copy the suffix
\texttt{\_raw}, so a name without the suffix is always trimmed. The joint units are part of the entry because the two
cube-in-box recordings use LeRobot's normalised range and the sorting recording uses degrees; datasets in different
units are never mixed, and errors are only compared within one dataset.

\chart{fig_datasets}{The datasets: episodes and episode lengths, before and after trimming.}{fig:datastats}

\section{Summary}
\textslot{the chapter's summary: the tasks, the recording, trimming and the versioned datasets}
